\chapter{Conclusion}

For measurement of air method 2 yielded $n = 1.0002103 \pm 5.3 \cdot 10^{-6}$. 
While this value is lower than the accepted literature value of $n \approx 1.000277$, 
the discrepancy can be attributed to systematic errors in the manual pressure readings. 
An analysis of the uncertainty suggests that a deviation in the starting gauge pressure of approximately $0.1$ to $0.14 \text{ atm}$ which is well within the range of human visual error during a dynamic pressure drop.
This would account for the difference, bringing the results into alignment with expected values.


The second part of the experiment characterized the effect of a phase retarder on interference intensity. 
The derived theoretical model,
$$I = I_0 (1 + \cos\Delta \phi \cos(2\theta))$$
predicts that the interference visibility is modulated by the orientation of the HWP's fast axis relative to the incident polarization.
By using a piezo-mounted mirror to sweep through $\Delta \phi$, 
we were able to isolate the maximum intensities and fit the data to the model $I \propto 1 + |\cos(2\theta)|$. 